\chapter{Your First Session}
\label{ch:first-session}

This chapter is a guided tour. Work through it at the Breadbox prompt; each exercise introduces one idea and gives you a way to check what happened.

\section{Meet the prompt}
After boot and login, try:

\begin{lstlisting}
help
pwd
ls /
cat /etc/motd
\end{lstlisting}

The prompt is where commands are entered. \cmd{help} displays the help text installed on the current disk; \cmd{pwd} reports the working directory; \cmd{ls} lists directory entries; and \cmd{cat} writes a file to the screen.

\section{Make a place to work}
\begin{lstlisting}
mkdir /home/work
cd /home/work
pwd
\end{lstlisting}

\section{Create and inspect a file}
\begin{lstlisting}
echo hello > greeting
cat greeting
echo again >> greeting
cat greeting
\end{lstlisting}

The first redirection creates or replaces the file. The second appends to it.

\section{Connect programs}
\begin{lstlisting}
cat greeting | wc
\end{lstlisting}

The pipe connects the output of \cmd{cat} to the input of \cmd{wc}. This is not a special built-in text operation: it connects two processes.

\begin{trythis}[Experiment]
Replace \cmd{cat} with \cmd{echo}, then try piping the output through \cmd{grep}. Change the search word and observe the result.
\end{trythis}

\section{Start a background job}
\begin{lstlisting}
count &
jobs
ps
\end{lstlisting}

The ampersand asks the shell to run the command in the background. The job list and process table show different views of running work. Use RUN/STOP to interrupt the foreground command; use \cmd{kill} when you need to signal a process by PID.
